\documentclass[10pt,a4paper]{amsart}
\usepackage[margin=19mm]{geometry}
\usepackage{amsmath,amssymb,mathtools}
\usepackage[scr=rsfs,cal=euler]{mathalfa}
\usepackage{xfrac}
\usepackage[unicode,hidelinks,bookmarks=false]{hyperref}
\newcommand{\ie}{i.e.\ }
\newcommand{\cf}{cf.\ }
\newcommand{\resp}{respectively\ }
\newcommand{\Subsection}{\subsection}
\providecommand{\nobreakdash}{\nobreak\hbox{-}}
\setlength{\emergencystretch}{2em}
\multlinegap0pt
\usepackage{bm}
\usepackage{bbm}
\usepackage{dsfont}
\usepackage{xspace}
\usepackage{cancel}
\usepackage{mathtools}
\usepackage{amsthm}
\makeatletter
\@ifundefined{theo}{\newtheorem{theo}{Theo}}{}
\@ifundefined{lemm}{\newtheorem{lemm}[theo]{Lemma}}{}
\makeatother
\DeclareMathOperator{\Spec}{Spec}
\title[Magic angle (in)stability and mobility edges in disordered C]{Magic angle (in)stability and mobility edges in disordered Chern insulators}
\author{Simon Becker, Izak Oltman, Martin Vogel}
\date{Source: arXiv:2309.02701v2 (CC BY 4.0)}
\begin{document}
\maketitle

\noindent\emph{Source and licence.} Magic angle (in)stability and mobility edges in disordered Chern insulators. Version arXiv:2309.02701v2 (CC BY 4.0), distributed under the Creative Commons Attribution 4.0 International licence, as recorded in the arXiv licence metadata for this exact version and shown on the abstract page https://arxiv.org/abs/2309.02701v2. Full rights notice: licenses/arxiv-2309.02701-CC-BY-4.0.txt.\par\smallskip

\noindent\textit{Editorial scope.} Counted unit: the Lemm whose source label is \texttt{lemm:SLIEDI} in the article (source0.tex lines 1157--1166, source label \texttt{lemm:SLIEDI}), delivered together with the article's own complete original proof of it (source0.tex lines 1167--1189, one \texttt{proof} environment of 23 lines). No mathematics is written, repaired or re-derived: statement and proof are the article's own text line for line. Required context retained uncounted: none. Every definition, display and label of those lines is retained unchanged. Omitted from this package and not redistributed: 1--1156, 1190--1786. Nothing inside a retained excerpt is omitted or altered: every retained body is delivered exactly as the article's lines are written, so no omission receipt of any kind accompanies this package. Citations: the retained excerpts carry no citation command at all, so no bibliography entry of the article is transcribed and no reference is redistributed. Reference adaptation: none was needed and none was made; every reference inside the delivered unit resolves inside it, so no unresolved reference or citation is produced. Printed slips kept literal: no printed word, symbol, hypothesis or number of the counted statement or of its proof was altered. Layout and class: the article's own class is \texttt{amsart} with packages including amsmath, amssymb, amsthm, graphicx, mathrsfs, url, color, hyperref, amsxtra, tikz, wasysym, subcaption, array, ulem; the wrapper is the lane's pinned \texttt{amsart} editorial wrapper at 10pt on A4, which restates the theorem-like environments the retained text needs and adds the article's own macro definitions that the retained text needs (\texttt{\textbackslash Spec}), with extra wrapper packages (bm, bbm, dsfont, xspace, cancel, mathtools). The delivered numbering is the unit's own. Assets: no retained span contains a figure environment, a graphics-inclusion command, a PDF-page inclusion, a table or a TikZ picture; no image file is copied into this package, the package declares no asset dependency and no image file is redistributed with it. Licence: version arXiv:2309.02701v2 of this article is distributed under the Creative Commons Attribution 4.0 International licence, as recorded in the arXiv licence metadata for this exact version and shown on the abstract page https://arxiv.org/abs/2309.02701v2. A full rights notice is delivered at licenses/arxiv-2309.02701-CC-BY-4.0.txt. Characters: every retained excerpt is pure ASCII, so the delivered engine, XeLaTeX with its default Latin Modern text font, reports no missing character.
\begin{lemm}[SLI \& EDI]
\label{lemm:SLIEDI}
Let $J$ be a compact interval. For $L, l',l'' \in 2\mathbb N$ and $x,y',y'' \in \Gamma$ with $\Lambda_{l''}(y) \subsetneq \Lambda_{l'}(y') \subsetneq \Lambda_L(x)$, then $\mathbf P$-almost surely: If $E \in J \cap (\Spec(H_{\lambda,\Lambda_L(x)}) \cap \Spec(H_{\lambda, \Lambda_{l'}(y')}))^c$ then the Simon-Lieb inequality holds
\[\tag{SLI} \begin{split} \Vert \Xi_{\Lambda_L(x)}(H_{\lambda,\Lambda_L(x)}-E)^{-1} \chi_{\Lambda_{l''}(y)}\Vert &\lesssim_J \Vert \Xi_{\Lambda_{l'}(y')}(H_{\lambda,\Lambda_{l'}(y')}-E)^{-1} \chi_{\Lambda_{l''}(y)} \Vert \\
&\times \Vert \Xi_{\Lambda_L(x)} (H_{\lambda,\Lambda_L(x)}-E)^{-1} \Xi_{\Lambda_{l'}(y')}\Vert.\end{split}\]
Moreover, for any $z \in \Gamma$ and any generalized eigenfunction $\psi$\footnote{ $\psi$ solving $(H_{\lambda}-E)\psi=0$ and growing at most polynomially} with generalized eigenvalue $E \in J \cap \Spec(H_{\lambda,\Lambda_L(x)})^c$ satisfies $\mathbf P$-almost surely the eigenfunction decay inequality
 \[\tag{EDI}
\begin{split} \Vert \chi_z \psi \Vert &\lesssim_J \Vert \Xi_{\Lambda_L(x)} (H_{\lambda,\Lambda_L(x)}-E)^{-1}\chi_z \Vert \Vert \Xi_{\Lambda_L(x)} \psi \Vert.
\end{split}\]
\end{lemm}

\begin{proof}
\begin{enumerate}
\item The proof of the SLI can be streamlined for linear differential operators with disorder of Anderson-type. We start from the following resolvent identity
\[ \begin{split}
(H_{\lambda}-E)\tilde \chi_{\Lambda_{l'}(y')}(H_{\lambda,\Lambda_L(x)}-E)^{-1} &= [ (H_{\lambda}-E),\tilde \chi_{\Lambda_{l'}(y')}](H_{\lambda,\Lambda_L(x)}-E)^{-1} \\
&\ + \tilde\chi_{\Lambda_{l'}(y')} (H_{\lambda}-E) (H_{\lambda,\Lambda_L(x)}-E)^{-1}.
\end{split}\]
Using that by assumption $\Lambda_{l'}(y') \subset \Lambda_{L}(x)$ we have $\chi_{\Lambda_{l'}(y')}H_{\lambda} = \chi_{\Lambda_{l'}(y')}H_{\lambda,\Lambda_L(x)}$ and find by substituting $\chi_{\Lambda_{l'}(y')}H_{\lambda} $ in the last line above 
\[ (H_{\lambda}-E)\tilde \chi_{\Lambda_{l'}(y')}(H_{\lambda,\Lambda_L(x)}-E)^{-1} = [ H_{\lambda},\tilde \chi_{\Lambda_{l'}(y')}](H_{\lambda,\Lambda_L(x)}-E)^{-1} + \tilde \chi_{\Lambda_{l'}(y')}.\]
Since $H_{\lambda}  \tilde\chi_{\Lambda_{l'}(y')} = H_{\lambda,\Lambda_{l'}(y')} \tilde\chi_{\Lambda_{l'}(y')}$ we find by multiplying the previous line by $(H_{\lambda,\Lambda_{l'}(y')}-E)^{-1}$ that 
\[ \begin{split} \tilde \chi_{\Lambda_{l'}(y')}(H_{\lambda,\Lambda_L(x)}-E)^{-1} &= (H_{\lambda,\Lambda_{l'}(y')}-E)^{-1}[ H_{\lambda,\Lambda_{l'}(y'),},\tilde \chi_{\Lambda_{l'}(y')}](H_{\lambda,\Lambda_L(x)}-E)^{-1} \\
&\quad + (H_{\lambda,\Lambda_{l'}(y')}-E)^{-1}\tilde\chi_{\Lambda_{l'}(y')}.
\end{split}\] 
Multiplying this equation from the left by $\chi_{\Lambda_{l''}(y)}$ and from the right by $\Xi_{\Lambda_L(x)}$, the SLI ready follows from the boundedness of $[ H_{\lambda,\Lambda_{l'}(y'),},\tilde \chi_{\Lambda_{l'}(y')}]$ and submultiplicativity of the operator norm, as $\tilde \chi_{\Lambda_{l'}(y')} \Xi_{\Lambda_L(x)}=0$ implies that the second term on the right vanishes and 
\begin{equation}
\label{eq:commutator}
[ H_{\lambda,\Lambda_{l'}(y'),},\tilde \chi_{\Lambda_{l'}(y')}] = \Xi_{\Lambda_{l'}(y')}[ H_{\lambda,\Lambda_{l'}(y'),},\tilde \chi_{\Lambda_{l'}(y')}] \Xi_{\Lambda_{l'}(y')}.
\end{equation}

\item For the proof of the EDI, it suffices to choose $\psi$ as in the Lemma and observe the resolvent identity $(H_{\lambda, x, L}-E)^{-1}[H_{\lambda}, \tilde \chi_{\Lambda_L(x)} ]\psi=\tilde \chi_{\Lambda_L(x)} \psi$ which is easily verified by using $(V_X-V_{X,\Lambda_L(x)}) \tilde \chi_{\Lambda_L(x)}=0$ and $H_{\lambda} \psi = E \psi.$
Using then an analog of \eqref{eq:commutator}, $[H_{\lambda}, \tilde \chi_{\Lambda_L(x)} ]= \Xi_{\Lambda_L(x)}[H_{\lambda}, \tilde \chi_{\Lambda_L(x)} ]\Xi_{\Lambda_L(x)},$ together with the boundedness of the commutator shows the claim.
\end{enumerate}
\end{proof}

\end{document}
